\chapter{CONCLUSIONES}

\section{Conclusiones del trabajo}

El alcance del proyecto: diseñar, desarrollar y validar un sistema escalable de microservicios para la ingesta, el procesamiento y el análisis de datos IoT, aplicado a la supervisión en tiempo real de la telemetría de consumo energético, fue cumplido en su totalidad.

Se probó la madurez y capacidad de las herramientas de código abierto para implementar un sistema de tratamiento de datos en flujo, resiliente y escalable, presentándolo como alternativa ante la amplia oferta de herramientas gestionadas y sistemas en la nube.

Con la debida consideración de las limitaciones en este proyecto, se lograron cumplir los diversos objetivos a lo largo de su desarrollo:

\begin{enumerate}
  \item \textbf{Persistencia de telemetría rápida y fiable.} El 95~\% de la telemetría validada se persistió en un tiempo inferior a los 25~segundos, con 0~\% de mensajes perdidos durante el procesamiento.
  \item \textbf{Gobernanza del esquema de datos.} El 100~\% de los eventos ingeridos se validaron contra el esquema Avro registrado, sin ningún fallo de deserialización ni mensajes erróneamente desviados al tópico de inválidos.
  \item \textbf{Procesamiento en flujo con baja latencia.} Los micro-lotes de Spark Structured Streaming se procesaron por debajo de los 2~segundos, a una tasa estable sin incremento de tiempo.
  \item \textbf{Visualización diferenciada operacional y analítica.} Se diseñaron e implementaron los dashboards de Grafana y las páginas de Power BI con las visualizaciones del consumo energético, la eficiencia operativa, la detección de anomalías y el análisis histórico.
  \item \textbf{Resiliencia del sistema.} Se comprobó la capacidad de auto recuperación de los componentes del sistema en menos de 60~segundos con 0~\% de pérdida de datos.
\end{enumerate}

Finalmente, la containerización, documentación y dispocisión del proyecto al público, lo hace disponible para su adopción, mejoramiento y adaptación a otros proyectos similares de monitorización energética.
